% cuadro C13_4
\begin{table}[htbp]
\centering
\scriptsize
\caption{Descomposición del cambio de la razón SHRFSP a PIB en cuatro variables, con la compensación por inflación como rango}
\label{cua:C13_4}
\begin{tabular}{p{3.6cm}rrrrrrr}
\toprule
\textbf{Convención del pago} & \textbf{RFSP} & \textbf{Crecim.} & \textbf{Inflación den.} & \textbf{Tipo cambio} & \textbf{Suma} & \textbf{Compens.} & \textbf{Infl. neta} \\
\midrule
deflactor del PIB, 4.8 & 4,10 & -1,10 & -2,40 & -0,61 & -0,01 & 2,34 & -0,05 \\
INPC promedio, 3.5 & 4,10 & -1,10 & -2,40 & -0,61 & -0,01 & 1,71 & -0,69 \\
INPC dic/dic, 3.0 & 4,10 & -1,10 & -2,40 & -0,61 & -0,01 & 1,46 & -0,93 \\
solo lo identificable: adecuaciones, 0.3 \% del PIB & 4,10 & -1,10 & -2,40 & -0,61 & -0,01 & 0,30 & -2,10 \\
\bottomrule
\end{tabular}
\\[2pt]\begin{minipage}{0.92\textwidth}\scriptsize Puntos del PIB. Las cuatro primeras columnas no dependen de la convención y suman el cambio del acervo. La compensación por inflación pagada NO es un término aditivo del marco: es una lectura del RFSP, y el paquete no identifica qué índice le corresponde. Se presenta como rango con su cota, nunca como punto. Elaborado por el ITED con datos del CGPE 2026, Anexos III.1 y III.2.\end{minipage}
\end{table}
